\documentclass[reqno, 12pt]{amsart}

\usepackage{amssymb}
\usepackage{amsthm}
\usepackage{amsmath}
\usepackage{amsxtra}
\usepackage{mathrsfs}
\usepackage{comment}
\usepackage{graphicx}
\usepackage{placeins}
\usepackage{multicol}
\usepackage{bbm}
\usepackage{stmaryrd}
%\usepackage{enumerate}
\usepackage[inline, shortlabels]{enumitem}
\usepackage{mathdots}
\usepackage{colonequals}
\usepackage{hyperref}


\usepackage{calc}
\usepackage{tikz}
\usetikzlibrary{arrows,calc,automata,shadows,backgrounds,positioning,intersections,fadings,decorations.pathreplacing,shapes,matrix}
\usepackage{tikz-cd}

\usepackage{fullpage}

\newtheorem{thm}{Theorem}[section]
\newtheorem{lem}[thm]{Lemma}
\newtheorem{cor}[thm]{Corollary}
\newtheorem{conj}[thm]{Conjecture}
\newtheorem{prop}[thm]{Proposition}
\theoremstyle{definition}  
\newtheorem{defn} [thm] {Definition} 
\newtheorem{claim}[thm]{Claim}
\newtheorem{example} [thm] {Example}
\newtheorem{rem} [thm] {Remark}

\newenvironment{problab}[1]
{\noindent\textbf{Problem #1}.}
{\vskip 6pt}

\newenvironment{exolab}[1]
{\noindent\textbf{Exercise #1}.}
{\vskip 6pt}

\theoremstyle{remark}
\newtheorem*{soln}{Solution}

\newcommand{\C}{\mathbb C}
\renewcommand{\H}{\mathbb H}
\newcommand{\D}{\mathbb D}
\newcommand{\F}{\mathbb F}
\newcommand{\Q}{\mathbb Q}
\newcommand{\R}{\mathbb R}
\newcommand{\Z}{\mathbb Z}
\newcommand{\T}{\mathbb T}
\renewcommand{\P}{\mathcal P}
\renewcommand{\O}{\mathcal O}
\newcommand{\m}{\mathfrak m}
\newcommand{\A}{\mathbb A}
\renewcommand{\SS}{\mathbb S}
\renewcommand{\L}{\mathcal L}

\newcommand{\B}{\mathcal B}
\newcommand{\E}{\mathcal E}

\newcommand{\SSpan}{\text{span}}

\DeclareMathOperator{\lcm}{lcm}
\DeclareMathOperator{\img}{img}
\DeclareMathOperator{\Ann}{Ann}
\DeclareMathOperator{\Tor}{Tor}
\DeclareMathOperator{\tr}{tr}
\DeclareMathOperator{\Hom}{Hom}
\DeclareMathOperator{\End}{End}
\DeclareMathOperator{\Aut}{Aut}
\DeclareMathOperator{\Gal}{Gal}
\DeclareMathOperator{\sgn}{sgn}
\DeclareMathOperator{\ord}{ord}
\DeclareMathOperator{\ad}{ad}
\DeclareMathOperator{\coker}{coker}
\DeclareMathOperator{\rank}{rank}
\DeclareMathOperator{\minpoly}{minpoly}

\DeclareMathOperator{\id}{id}

\usepackage{mathpazo}

\everymath{\displaystyle}

\tikzset{commutative diagrams/.cd, arrow style = tikz, diagrams = {>=latex}}

\newenvironment{amatrix}[1]{%
  \left(\begin{array}{@{}*{#1}{c}|c@{}}
}{%
  \end{array}\right)
}

\usepackage{abstract}
\renewcommand{\abstractname}{}    % clear the title
\renewcommand{\absnamepos}{empty} % originally center
%\setlength{\abstitleskip}

\begin{document}
%\renewcommand{\abstractname}{\vspace{-\baselineskip}}

\title{18.700 Problem Set 6}
%\author{Évariste Galois}
\dedicatory{Due Wednesday, October 30 at 11:59 pm on \href{https://canvas.mit.edu/courses/27315}{Canvas}}
%\thanks{}
%\contrib{}
%\address{}

\maketitle

\begin{abstract}
  \noindent Collaborated with:\\
  Sources used: 
\end{abstract}
\vskip 0.5cm

%\begin{problab}{1} (5B \#12) (10 points)
%  Define $T \in \L(\F^n)$ by
%  \begin{equation*}
%    T(x_1, x_2, x_3, \ldots, x_n) = (x_1, 2 x_2, 3 x_3, \ldots, n x_n) \, .
%  \end{equation*}
%  Find the minimal polynomial of $T$.
%\end{problab}

\begin{problab}{1} (5B \#21) (5 points)
  Suppose $V$ is finite-dimensional and $T \in \L(V)$. Prove that the minimal polynomial of $T$ has degree at most $1 + \dim(\img(T))$. (Observe that if $\dim(\img(T)) < \dim(V) - 1$, then this exercise gives a better upper bound for the degree of $\minpoly(T)$.) (\textit{Hint}: Recall that $\img(T)$ is $T$-invariant.)
\end{problab}

%\begin{soln}
%\end{soln}

\begin{problab}{2} (5B \#22) (6 points)
  Suppose $V$ is finite-dimensional and $T \in \L(V)$. Prove that $T$ is invertible iff $I \in \SSpan(T, T^2, \ldots, T^{\dim(V)})$.
\end{problab}

\begin{problab}{3} (5C \#6) (5 points)
  Suppose $\F = \C$, $V$ is finite-dimensional and $T \in \L(V)$. Prove that $V$ has a $k$-dimensional $T$-invariant subspace for each $k = 1, \ldots, \dim(V)$.
\end{problab}

\begin{problab}{4} (5C \#11) (6 points)
  A square matrix is \textit{lower-triangular} if all its entries above the diagonal are $0$. Suppose $\F = \C$ and $V$ is finite-dimensional. Prove that if $T \in \L(V)$, then there exists a basis of $V$ with respect to which $T$ has a lower-triangular matrix.
\end{problab}

\begin{problab}{5} (5D \#1) (7 points)
  Suppose $V$ is a finite-dimensional $\C$-vector space and $T \in \L(V)$.
  \begin{enumerate}[(a)]
    \item 
      Prove that if $T^4 = I$, then $T$ is diagonalizable.
    \item
      Prove that if $T^4 = T$, then $T$ is diagonalizable.
    \item
      Give an example of an operator $T \in \L(\C^2)$ such that $T^4 = T^2$ and $T$ is not diagonalizable. (Make sure to justify why your $T$ is not diagonalizable.)
  \end{enumerate}
\end{problab}

\begin{problab}{6} (5D \#21) (12 points)
  The \textit{Fibonacci sequence} $F_0, F_1, F_2, \ldots$ is defined by
  \begin{equation*}
    F_0 = 0, \ F_1 = 1, \ \text{and} \ F_n = F_{n-2} + F_{n-1} \ \text{for $n \geq 2$.}
  \end{equation*}
    Define $T \in \L(\R^2)$ by $T(x,y) = (y, x+y)$.
    \begin{enumerate}[(a)]
      \item 
        Show that $T^n(0,1) = (F_n, F_{n+1})$ for each $n \in \Z_{\geq 0}$.
      \item
        Find the eigenvalues of $T$.
      \item
        Find a basis of $\R^2$ consisting of eigenvectors of $T$.
      \item
        Use the solution to (c) to compute $T^n(0,1)$. Conclude that
        \begin{equation*}
          F_n = \frac{1}{\sqrt{5}} \left[ \left( \frac{1+\sqrt{5}}{2} \right)^n - \left( \frac{1 - \sqrt{5}}{2} \right)^n \right]
        \end{equation*}
        for each $n \in \Z_{\geq 0}$. (The number $(1+\sqrt{5})/2$ is called the \textit{golden ratio}.)
    \end{enumerate}
\end{problab}

\end{document} 	
